\documentclass[11pt]{article}
\usepackage[a4paper,margin=1in]{geometry}
\usepackage{amsmath,amssymb,amsthm}
\usepackage{mathtools}
\usepackage{booktabs}
\usepackage{graphicx}
\usepackage{xcolor}
\usepackage{hyperref}
\hypersetup{colorlinks=true,linkcolor=blue,urlcolor=blue,citecolor=blue}
\newcommand{\rt}[1]{\texttt{#1}}
\newcommand{\fn}[1]{\texttt{#1}}
\newcommand{\K}{\mathcal{K}}
\newcommand{\ifgraphic}[2]{\IfFileExists{#1}{\includegraphics[width=#2]{#1}}%
  {\fbox{\parbox{#2}{\centering\small figure \texttt{\detokenize{#1}} pending}}}}
\newcommand{\iftable}[1]{\IfFileExists{#1}{\input{#1}}%
  {\fbox{\parbox{0.9\textwidth}{\centering\small table
  \texttt{\detokenize{#1}} pending}}}}

\theoremstyle{plain}
\newtheorem{theorem}{Theorem}
\newtheorem{proposition}{Proposition}
\theoremstyle{remark}
\newtheorem{remark}{Remark}

\title{\textbf{Report R31 --- the same four rules with $V=X$:\\
what survives of R30 when the gate becomes a permutation}}
\author{Tier 1 analysis}
\date{2026-08-28}

\begin{document}
\maketitle

\begin{abstract}
R30 derived the Krylov sector-size distribution of $W_{156}/W_{198}$ and
$W_{108}/W_{201}$ with a Hadamard gate. Replacing the gate by $X$ makes the
circuit a \emph{permutation}, so the sectors become its orbits, and the question
is how much of the R30 machinery is gate-independent. The answer is sharp on both
sides. What dies is one thing: R30 rests on the flip reduction, which needs
$\langle b|V|b\rangle\ne0$ so that a half-layer image is a full \emph{subcube};
$X$ has a zero diagonal, the choice at an active site is forced, and the
transitivity half of R30's Theorem~1 goes with it. What survives is more than one
would guess. The wall set is still conserved --- for free, since
$\mathrm{supp}(U_X|x\rangle)\subseteq\mathrm{supp}(U_H|x\rangle)$ --- so the label
is still an invariant, merely no longer a complete one, and the $X$ orbits
strictly refine the Hadamard sectors. And, measured rather than assumed, inside a
label fibre the $X$ map is the \textbf{direct product} of independent
per-segment permutations, each depending only on the segment. So R30's renewal
factorisation survives one level down, with its scalar gap kernel $K$ replaced by
a \textbf{cycle type} $C$ satisfying $\sum C=K$, and its product replaced by an
\textbf{lcm}: choosing one cycle $c_j$ per segment yields $\prod_j c_j/\mathrm{lcm}_j c_j$
orbits of length $\mathrm{lcm}_j c_j$. For the ascent pair $C$ is a single cycle,
so the orbit length is literally the lcm of R30's kernels; for the symmetric pair
$C$ is the spectrum of the hard-core $X$ automaton on a segment, whose longest
cycle is $3\ell-7$. This reproduces the exhaustive orbit multiset term by term on
all $80$ (rule, boundary, $N$) triples with $N\le13$. The consequences are
physical: replacing a product by an lcm turns the largest sector from
$\varphi^N$ into a \textbf{Landau-type maximum}, $\exp\Theta(\sqrt{N\log N})$,
which we identify exactly for the ascent pair --- so under $X$ there is no
macroscopic sector at all, $D_{\max}$ is not even monotone in $N$, and R30's
log-normal law has no analogue.
\end{abstract}

\tableofcontents

\section{The question}

The four rules and their wall words are those of R30:
$W_{201}=\rt{VIII}$ (word \rt{11}), $W_{108}=\rt{IIIV}$ (\rt{00}),
$W_{156}=\rt{IIVI}$ (\rt{01}), $W_{198}=\rt{IVII}$ (\rt{10}). Everything below
replaces the Hadamard by $X$ and keeps every other convention fixed: even sites
then odd, sequential within a layer, \rt{obc0} $=$ frozen $0$ pads, and sectors
$=$ weakly connected components of the one-cycle graph. Both boundary conditions
are carried throughout, and all four rules were computed from their own wall
words --- the coincidences below are results, not inputs.

\section{What breaks: exactly one thing}

\begin{proposition}[the load-bearing use of the Hadamard]
R30 reduces the sector partition to the connected components of the single-flip
graph. The proof needs both matrix elements of the gate to be non-zero, so that
the image of a half-layer from a basis state is the full subcube spanned by the
active sites. For $V=X$ the diagonal element vanishes: ``do not flip'' has
amplitude zero, the choice at an active site is forced, and the image is a single
point. The one-cycle map is a permutation of the $2^N$ basis states.
\end{proposition}

\noindent Two corollaries, both verified exhaustively for all four rules at both
boundary conditions:

\begin{itemize}
\item \textbf{Sectors are orbits.} There are no transients, every state is
recurrent, and the sum rule $\sum_\ell \ell\, n_N(\ell)=2^N$ is automatic rather
than a check.
\item \textbf{R30 Theorem~1's transitivity half dies.} Its proof moves one domain
wall by one site, or removes one particle at a time; those are single-site moves,
and under $X$ every active site flips at once. Nothing replaces them.
\end{itemize}

\section{What survives, and why}

\subsection{The label is still conserved}

\begin{proposition}
The wall set is constant on $X$ orbits, for all four rules and both boundary
conditions.
\end{proposition}

\begin{proof}
$\mathrm{supp}(U_X|x\rangle)\subseteq\mathrm{supp}(U_H|x\rangle)$, because the
$X$ image is one corner of the Hadamard subcube. Any quantity constant on
Hadamard sectors is therefore constant on $X$ orbits, and R30 Theorem~1 gives the
wall set as such a quantity.
\end{proof}

\noindent So the label is an invariant but not a complete one: the $X$ orbits
\emph{strictly refine} the Hadamard sectors. That is R11's refinement theorem,
now visible in the census --- at \rt{obc0}, $N=12$: $1442$ orbits inside $1081$
fibres for $W_{108}$, $1018$ inside $616$ for $W_{201}$, $688$ inside $377$ for
$W_{156}$.

A pleasant corollary in the other direction: a state with no active site is fixed
by any gate, so the \textbf{frozen sectors are gate-independent}. The fixed points
of the $X$ map are exactly R30's size-$1$ sectors, checked for $N=4\ldots12$, all
four rules, both boundary conditions. R30's closed forms for $n_{\rm frozen}$ ---
$\lfloor(N-1)/2\rfloor+2$ and $4/2$ for the ascent pair, $\Theta(\varphi^N)$ for
the symmetric one --- carry over unchanged.

\subsection{The renewal factorisation survives one level down}

This is the central measurement of the report.

\begin{theorem}[measured; \fn{check\_factorisation}]
\label{thm:factor}
Inside a label fibre the $X$ map is the \textbf{direct product} of independent
per-segment permutations, and each factor depends only on its own segment.
Zero failures and zero conflicts over every label at $N=10,12$, all four rules,
both boundary conditions; Tier-2 R-T21 reports the same over $N=3\ldots13$ from
an independent implementation.
\end{theorem}

\noindent The walls block information flow, so this is what one hopes for; but it
is not automatic, and it is what makes everything else possible. Given it, define
$C$, the \textbf{cycle type} of a segment, indexed by the segment's number of
states --- that is, by R30's kernel value. Then

\begin{equation}
\boxed{
\begin{aligned}
\text{R30 ($H$):}\quad & |\K| = L(g_0)K(g_1)\cdots R(g_k) && \text{a PRODUCT},\\
\text{R31 ($X$):}\quad & \text{choose } c_j\in C(g_j)\ \text{in each gap};\ \
\text{this gives } \tfrac{\prod_j c_j}{\mathrm{lcm}_j c_j} \ \text{orbits of
length } \mathrm{lcm}_j c_j .
\end{aligned}}
\end{equation}

\noindent Since $\sum C(g)=K(g)$, forgetting the cycle structure recovers R30
exactly: the cycle type refines the Hadamard kernel and nothing is lost.

\begin{center}
\small
\iftable{tab_r31_kernel.tex}
\end{center}

\paragraph{The ascent pair is trivial and the symmetric pair is not.}
For $W_{156}/W_{198}$, $C$ is a \emph{single} cycle of length equal to R30's
kernel value: the domain wall marches ballistically and reflects, visiting every
admissible position once per period. So the orbit length is literally the lcm of
R30's gap kernels, and a fibre of R30 size $\prod_j K_j$ carries
$\prod_j K_j/\mathrm{lcm}_j K_j$ orbits of length $\mathrm{lcm}_j K_j$. For
$W_{108}/W_{201}$ the segment carries the hard-core (PXP-like) $X$ automaton with
frozen ends, $|C(\ell)|=F(\ell+2)$ as in R30, and the spectrum is genuinely
non-trivial: $(2,3)$, $(3,5)$, $(2,3,8)$, $(3,7,11)$, $(2,3,5,10,14),\ldots$,
with the cycle lengths lying in arithmetic progressions of step $4$ and the
longest equal to $3\ell-7$ (verified here to $\ell=16$ and by Tier-2 R-T21 to
$\ell=26$).

\begin{remark}[indexing, so the two reports can be read together]
$C$ is indexed here by the pair (\emph{family}, \emph{number of states}), which
makes the corner class redundant --- keyed that way there are zero conflicts for
all four rules at both boundaries. The family is genuinely part of the key: at
every Fibonacci state count from $5$ up the two families carry different types
($v=5$ is $(5)$ for the ascent pair and $(2,3)$ for the symmetric one). Tier-2 R-T21 indexes by the geometric gap and therefore
needs a left/interior/right class; the two are the same table. Their
$C(L,\ell)=C(K,\ell+1)$ for $W_{201}$ is our ``both segments have $F(\cdot)$
states'', and their $C(L,\ell)=(1,)$ for $W_{156}$ is our ``a one-state segment
is a fixed point''. Their $\ell$ counts the collars and ours does not, so their
hyperbolic table is shifted by two.
\end{remark}

\subsection{The wall-free label: a segment on the chain, not on the ring}

R30's Remark~2 recorded that every correction in that report was the same term,
the wall-free label, and predicted a fourth. It happened twice more here, and the
truth turns out to be sharper than either guess that preceded it:

\begin{proposition}
At \rt{obc0} the wall-free label \emph{is} a segment --- its cycle type is exactly
$C$ evaluated at its state count, for all four rules, $N=4\ldots12$. At \rt{pbc}
it is not, for any of them.
\end{proposition}

\noindent The mechanism is the pad. On the chain the frozen $0$ outside each end
plays the part that a wall's collar plays in the bulk, so a wall-free chain is
simply one long segment. On the ring there is no boundary at all: for the ascent
rules the wall-free label is the two frozen uniform states, cycle type $(1,1)$ and
not the $(2,)$ of a two-state segment; for the symmetric rules it is the hard-core
$X$ automaton on a \emph{ring} of $L(N)$ states, whose spectrum
($(4)$, $(2,2,3)$, $(4,7)$, $(3,5,5,5)$, $(8,10,11)$ at $N=3\ldots7$) appears
nowhere among the segment kernels --- indeed usually \emph{no} segment has $L(N)$
states, since R30 showed Lucas numbers are generally not Fibonacci products.

\begin{remark}
Both of the errors here were mine. The first kernel extraction classified the
wall-free ring as a segment and duly reported the \rt{pbc} ascent kernel as
ambiguous at eight lengths; the correction to that then over-shot into ``the
wall-free label is never a segment'', which is false on the chain. Tier-2 R-T21
caught the second. The predictor was unaffected only because it computes the
wall-free type directly instead of looking it up.
\end{remark}

\subsection{The ring defect stops being a defect}

Under the Hadamard, $0^N$ and $1^N$ shared the empty label on the ring and had to
be separated by hand. Under $X$ they are simply two fixed points, and the lcm
assembly returns them with no special case at all.

\section{Verification}

\begin{figure}[htbp]
\centering
\ifgraphic{../../figures/fig_orbit_sizes.pdf}{\textwidth}
\caption{\textbf{The $X$-gate orbit-length distribution.}
(a) the measured histogram --- orbits of the $2^N$ permutation at $N=16$,
\rt{obc0} --- against the lcm assembly, overlaid as black dashes; they agree
term by term.
(b) the largest sector under $X$ (solid) against the Hadamard (dashed). The
symmetric pair falls off its $\varphi^N$ line immediately; the ascent pair
separates later, because its Hadamard $D_{\max}=4^{N/5}$ is itself slow. The $X$
curves are not monotone.
(c) the new kernel: for the symmetric pair the segment cycle spectrum, with the
longest cycle $3\ell-7$; for the ascent pair a single cycle of length $\ell+1$.}
\label{fig:orbits}
\end{figure}

The assembly reproduces the exhaustive orbit multiset of the $2^N$ permutation
\textbf{term by term} on all $80$ (rule, boundary, $N$) triples with
$N=4\ldots13$ --- not merely the counts and the maxima. Tier-2 R-T21 reports the
same on $96$ triples to $N=14$ from an independent implementation, and our orbit
counts agree to the digit.

\begin{center}
\small
\textbf{\rt{obc0}} --- $n_{\rm sectors}$ ($X$/$H$) and $D_{\max}$ ($X$/$H$)\\[2pt]
\resizebox{\textwidth}{!}{\iftable{tab_r31_census_obc0.tex}}\\[8pt]
\textbf{\rt{pbc}}\\[2pt]
\resizebox{\textwidth}{!}{\iftable{tab_r31_census_pbc.tex}}
\end{center}

\section{The distribution: an arithmetic law, not a log-normal}

\subsection{$D_{\max}$ is a Landau-type maximum}

Replacing a product by an lcm is what changes the physics. For the ascent pair
the orbit length is exactly $\mathrm{lcm}$ of R30's gap kernels, so

\begin{equation}
\boxed{\;
D_{\max} \;=\; \max\ \mathrm{lcm}(g_1,\ldots,g_k)
\quad\text{subject to}\quad
\sum_j (g_j+1)\ \begin{cases}\le N+1 & (\rt{obc0}),\\ = N & (\rt{pbc}).\end{cases}}
\end{equation}

\noindent Zero mismatches against exhaustive enumeration, $N=4\ldots19$, both
boundary conditions, and independently against Tier-2 R-T21's orbit multisets to
$N=18$.

\begin{remark}[an implementation trap]
Maximum-lcm has \emph{no optimal substructure}: the obvious memoised recursion
$\max_g\, g\cdot\mathrm{Landau}(\text{budget}-g-1)$ is wrong, because the best
sub-solution for a smaller budget need not be the one that combines best with
$g$. It returns $45$ instead of $60$ at $N=15$. The implementation here carries
the whole frontier of achievable lcms with their minimal costs, and agrees with
brute force; anyone reproducing this should check $N=15$ first.
\end{remark} (The ring needing the budget spent \emph{exactly} rather
than at most is the same cyclic-budget effect as R30's parity selection rule.)
This is Landau's function with a cost penalty, so
$D_{\max}=\exp\Theta(\sqrt{N\log N})$ --- \textbf{subexponential}. Measured,
$\ln D_{\max}/\sqrt{N\ln N}$ sits at $0.62,0.67,0.69,0.69,0.74$ for $W_{156}$ at
\rt{obc0}, $N=12,16,20,24,28$.

So the single most consequential difference between the two gates is this: under
the Hadamard the largest sector holds $\varphi^N$ states, an extensive fraction of
a fragmented space; under $X$ \emph{no sector is macroscopic at all}. And because
an lcm is arithmetic rather than monotone, $D_{\max}$ is not even increasing in
$N$ --- $W_{201}$ at \rt{pbc} gives $56$, $166$, $153$ at $N=16,17,18$. That is
not noise and not a bug; it is the signature that this is no longer a growth law.

\subsection{The sector count, and a convergence between the families}

The two families have almost the same growth base under $X$, which they did not
under the Hadamard:
\begin{center}
\small
\begin{tabular}{llrr}
\toprule
& & base under $X$ & base under $H$ \\
\midrule
\rt{obc0} & $W_{156}/W_{198}$ & $1.8299$ & $\varphi=1.6180$ \\
\rt{obc0} & $W_{108}/W_{201}$ & $1.8412$ & $\rho^2=1.7549$ \\
\rt{pbc}  & $W_{156}/W_{198}$ & $1.8124$ & $1.6164$ \\
\rt{pbc}  & $W_{108}/W_{201}$ & $1.8437$ & $1.7549$ \\
\bottomrule
\end{tabular}
\end{center}
(fitted over $N=10\ldots18$). More sectors, all of them smaller, and the
distinction between the Fibonacci and unipotent families --- which organised
every result in R30 --- largely washes out of the \emph{count}, while remaining
sharp in the kernel and in $D_{\max}$.

\subsection{What replaces the log-normal, and what we are not claiming}

R30's log-normal came from $\ln|\K|$ being a renewal-\emph{reward} sum: an
additive functional of a renewal chain, hence a central limit theorem. Here
$\ln(\text{orbit length})=\ln\mathrm{lcm}(c_1,\ldots,c_k)$, which is not additive
--- it is $\sum_p \max_j v_p(c_j)\ln p$, an arithmetic object. Empirically the
typical orbit is strongly sub-linear in $N$:
\begin{center}
\small
\begin{tabular}{lrrrrr}
\toprule
$N$ & $12$ & $16$ & $20$ & $24$ & $28$ \\
\midrule
$\langle\ln \ell\rangle$ ($W_{156}$, \rt{obc0}) & $1.5624$ & $1.8750$ & $2.1173$ & $2.3198$ & $2.4926$ \\
$\;/\ln N$    & $0.629$ & $0.676$ & $0.707$ & $0.730$ & $0.748$ \\
$\;/\ln^2 N$  & $0.253$ & $0.244$ & $0.236$ & $0.230$ & $0.224$ \\
\bottomrule
\end{tabular}
\end{center}
The first ratio rises and the second falls, and at these sizes the two candidate
laws are not separated. They have since been separated, by Tier-2 R-T21, and the
answer is $\Theta(\ln N)$.

\begin{proposition}[Tier-2 R-T21; anchor reproduced here]
\label{prop:logn}
Under the \emph{state} measure the mean log orbit length is
$\langle\ln\ell\rangle \sim \ln N/\ln 2$, and the orbit measure obeys the same
order.
\end{proposition}

\noindent The route is worth recording because it removes the enumeration
entirely. Since $\ln\mathrm{lcm}=\sum_p \ln p\,\max_j v_p(c_j)$ and the state
measure factorises over gaps, the probability that every gap has
$v_p(c_j)<e$ is a renewal sum with $K(g)$ replaced by the number of states of
that gap whose cycle length $c$ has $v_p(c)<e$; then
$\mathbb{E}[\max_j v_p]=\sum_{e\ge1}(1-\mathbb{P}(\cdot<e))$. For the ascent pair
each gap is a single cycle of length $g$, so that count is $g$ or $0$ and the
whole thing is closed form in $O(N^2)$. It gives, for $W_{156}$ at \rt{obc0},
\begin{center}
\small
\begin{tabular}{lrrrrrrr}
\toprule
$N$ & $28$ & $60$ & $120$ & $250$ & $500$ & $1000$ & $2000$\\
\midrule
$\langle\ln\ell\rangle_{\rm state}$ & $3.3138$ & $4.6269$ & $5.8388$ & $7.1116$ & $8.2812$ & $9.4119$ & $10.5925$\\
$\;/\ln N$   & $0.9945$ & $1.1301$ & $1.2196$ & $1.2880$ & $1.3325$ & $1.3625$ & $1.3936$\\
$\;/\ln^2 N$ & $0.2984$ & $0.2760$ & $0.2547$ & $0.2333$ & $0.2144$ & $0.1972$ & $0.1833$\\
\bottomrule
\end{tabular}
\end{center}
--- the first ratio climbing monotonically towards $1/\ln2=1.442695$ and the
second still falling at $N=2000$, which excludes $\ln^2N$. The $N=28$ entry is
$3.3138$, which is exactly the value the histogram gives here, so the two methods
agree at the one point where both can be evaluated.

The mechanism supplies the constant. Under the state measure the gaps form a
renewal sequence with $\mathbb{P}(g)\propto K(g)2^{-(g+1)}=g\,2^{-(g+1)}$, mean
block length $4$; so there are $\sim N/4$ gaps, the largest is $\sim\log_2 N$,
every integer below it occurs many times, and
$\ln\mathrm{lcm}\sim\psi(\text{largest gap})\sim\log_2 N$ with $\psi$ Chebyshev's
function. The approach is slow precisely because $\psi(G)/G\to1$ slowly, which is
the drift in the table.

Our measure is the orbit one, which is smaller --- the state measure is its
$\ell$-size-biased version and $\ln\ell$ is increasing, so
$\mathbb{E}_{\rm orbit}[\ln\ell]\le\mathbb{E}_{\rm state}[\ln\ell]$ by Chebyshev's
sum inequality. Hence the orbit mean is $O(\ln N)$ as well and $\ln^2N$ is
excluded for it too. The ratio of the two means drifts from $1.293$ at $N=12$ to
$1.319$ at $N=24$, which is why $N\le28$ could not separate the laws.

\paragraph{What is still open.} The \emph{constant} for the orbit measure. Its
weight is $\prod_j c_j/\mathrm{lcm}_j c_j$, and the lcm in the denominator does
not factorise over gaps the way $\sum_p\ln p\,\max_j v_p$ does, so the $O(N^2)$
route does not transfer. The order is settled; the prefactor is not. For the
symmetric pair the same argument gives $\Theta(\ln N)$ with a larger constant
(the gap law is $F(g-1)2^{-(g+1)}\sim(\varphi/2)^g$, so the largest gap is
$\sim\ln N/\ln(2/\varphi)$), but that is an argument and not a measurement,
because the exact route there needs the hyperbolic segment spectra and each costs
$\varphi^\ell$ states.

\section{Answering the question as asked}

\begin{enumerate}
\item \textbf{Is the old method still applicable?} Half of it, and the half that
survives is the interesting half. The wall grammar is gate-independent: the label
is still conserved, the fibre is still a product of free segments, and the
renewal decomposition still organises the calculation. What is gate-dependent is
the \emph{completeness} of the label, and with it every quantitative result of
R30, because those all computed the size of a label fibre and the fibre now
splits.

\item \textbf{Where did R30 explicitly need the Hadamard?} In exactly one place:
the flip reduction, which needs a non-zero diagonal so that a half-layer image is
a subcube. Everything else in R30 is downstream of that, through the transitivity
half of Theorem~1.

\item \textbf{What replaces it?} A cycle type in place of a scalar kernel, and an
lcm in place of a product. R30's $K(g)$ becomes $C(g)$ with $\sum C=K$, so the
Hadamard answer is recovered by forgetting the cycle structure --- the two reports
are the same combinatorics read at two resolutions.

\item \textbf{What changes physically?} The largest sector collapses from
$\varphi^N$ to $\exp\Theta(\sqrt{N\log N})$: under $X$ there is no macroscopic
Krylov sector, the fragmentation is finer everywhere, and the size distribution is
governed by lcm arithmetic rather than by a central limit theorem. Closed orbits
and permutations do change the distribution, and they change it in the direction
of \emph{more} fragmentation, not less.
\end{enumerate}

\section{Reproduction}

\begin{verbatim}
python -m qca_fragmentation.scaling.orbit_size_figure
python -m pytest qca_fragmentation/tests/test_orbit_sizes.py -q     # 66 tests
python -m pytest qca_fragmentation/tests -q                         # 2034 tests
\end{verbatim}

\noindent Engine in \fn{scaling/orbit\_sizes.py}, which reuses R30's
\fn{scaling/sector\_sizes.py} for the label decomposition, so the two reports
share one validated enumeration and differ only in the weight carried along it.
The complementary transfer-matrix treatment, the corner-class indexing, the
segment spectrum to $\ell=26$ and the ring wall-free spectra are Tier-2 report
R-T21.

\end{document}
